\chapter{Конкуренция выживших циклов в точном Gaussian-следе}
\label{ch:v7-weak-aligned-cycle-competition-gate}

\section{Полный одно-поколенческий носитель}

Рассмотрим полный $21$-мерный конечный оператор на вершинах
$H_{15}\oplus X_L\oplus X_R\oplus Y_L\oplus Y_R$. Singlet-фон содержит
три старых ребра и четыре ненулевые gauge-синглетные амплитуды. Тяжёлое
пространство двух выживших циклов имеет разложение
\begin{equation}
 \mathcal E_{
m heavy}^{\mathbb R}
 =\mathcal E_d^{\mathbb R}\oplus\mathcal E_W^{\mathbb R},
 \qquad \dim_{\mathbb R}\mathcal E_d=12,
 \quad \dim_{\mathbb R}\mathcal E_W=8.
 \label{eq:v7-weak-competition-heavy-split}
\end{equation}

Для чистого теплового профиля проверяется точный функционал
\begin{equation}
 S_t(\Phi)=\Tr e^{-t\Phi^2},\qquad t>0.
 \label{eq:v7-weak-competition-exact-gaussian}
\end{equation}

\section{Точный спектральный гессиан}

Если $\Phi_0=U\operatorname{diag}(\lambda_i)U^\dagger$ и
$E_a'=U^\dagger E_aU$, то формула Далецкого--Крейна даёт
\begin{equation}
 (H_t)_{ab}=\sum_{i,j}g_t(\lambda_i,\lambda_j)
 (E_a')_{ij}(E_b')_{ji},
 \qquad
 g_t(x,y)=\frac{f_t'(x)-f_t'(y)}{x-y},
 \quad f_t(x)=e^{-tx^2},
 \label{eq:v7-weak-competition-divided-hessian}
\end{equation}
с диагональным продолжением $g_t(x,x)=f_t''(x)$. Это вычисляет весь
экспоненциальный остаток без усечения по $a_6$.

Down- и weak-блоки гессиана ортогональны:
\begin{equation}
 H_t=H_d(t)\oplus H_W(t).
 \label{eq:v7-weak-competition-block-split}
\end{equation}
Однако их знаки различаются. При $t=1$ получено
\begin{align}
 \Spec H_d(1)&=\{0.041312^{\times6},\ 0.736618^{\times6}\},\nonumber\\
 \Spec H_W(1)&=\{-1.280223^{\times2},-1.004339,-0.558141,
 -0.265457^{\times2},-0.093662^{\times2}\}.
 \label{eq:v7-weak-competition-t-one-spectrum}
\end{align}
Следовательно, даже там, где цветная down-пара стабилизируется, все восемь
вещественных weak-направлений ещё отрицательны.

\section{Профильный проход и корневые вариации}

Логарифмический проход по $401$ точке
\begin{equation}
 10^{-4}\le t\le10^2
 \label{eq:v7-weak-competition-scan-domain}
\end{equation}
не нашёл ни одного $t$ с неотрицательным $H_W(t)$. При малом $t$ это
следует уже из $S_t=21-t\Tr\Phi^2+O(t^2)$; при большом $t$ отрицательный
вклад связи нулевой моды фонового пути $P_3$ с путём $P_6$ стремится к нулю
с отрицательной стороны.

Кроме того, singlet-фон не является стационарной точкой автономного
Gaussian-функционала. Например, при $t=1$ семь корневых производных равны
\begin{equation}
 (-4.414553,-4.414553,-0.549001,-0.635003,
  -0.866155,-0.093662,-1.407496),
 \label{eq:v7-weak-competition-root-tadpoles}
\end{equation}
и потому не исчезают.

\begin{theorem}[Запрет автономного точного Gaussian-родителя]
На полном $21$-мерном носителе точный функционал $\Tr e^{-t\Phi^2}$ не
продолжает singlet-Hodge вакуум в общий устойчивый родитель. Down-блок имеет
положительное окно, но weak-блок сохраняет отрицательную моду на всём
проверенном диапазоне и в обоих асимптотических пределах. Одновременно
ненулевые корневые производные сдвигают сам фон.
\end{theorem}

\begin{nogocriterion}[Запрет ручного сложения двух действий]
Нельзя исправить результат добавлением прежнего Hodge-потенциала с
произвольным весом: такой вес является новым параметром. Требуется один
точный функционал, чьи квадратичный, квартичный и циклический уровни
одновременно порождают singlet-вакуум и стабилизируют weak-конкурента.
\end{nogocriterion}

\begin{protocol}
\begin{enumerate}
 \item комплексная размерность полного носителя: \textbf{$21$};
 \item вещественная размерность тяжёлого сектора: \textbf{$20$};
 \item смешивание down--weak в гессиане: \textbf{ноль};
 \item положительное окно down-блока: \textbf{есть};
 \item положительное окно полного weak-блока: \textbf{нет};
 \item корневые tadpole автономного Gaussian: \textbf{ненулевые};
 \item точный Gaussian как общий родитель: \textbf{не проходит};
 \item итог: \textbf{no-go автономного профиля};
 \item следующий гейт: \textbf{единый точный Hodge--cycle функционал}.
\end{enumerate}
\end{protocol}

Машинный сертификат сохранён в
\path{s2t_v7_weak_aligned_cycle_competition_gate_results.json}.

\begin{thebibliography}{9}
\bibitem{v7-weak-daleckii}
V.~Noferini, \emph{A Daleckii--Krein Formula for the Fr\'echet Derivative
of a Generalized Matrix Function}, SIAM Journal on Matrix Analysis and
Applications 38 (2017), 434--457.
\bibitem{v7-weak-vassilevich}
D.~V.~Vassilevich, \emph{Heat Kernel Expansion: User's Manual}, Physics
Reports 388 (2003), 279--360; arXiv:hep-th/0306138.
\end{thebibliography}